From - Mon May 22 18:45:34 2000
Received: from grumpy.usu.edu (grumpy.usu.edu [129.123.1.86])
	by math.usu.edu (8.9.3+Sun/8.9.1) with ESMTP id LAA19951
	for <Symanzik@sunfs.math.usu.edu>; Sun, 21 May 2000 11:23:26 -0600 (MDT)
Received: from webmail.usu.edu ("port 2572"@webster.usu.edu [129.123.1.90])
 by cc.usu.edu (PMDF V5.2-32 #30472)
 with ESMTP id <01JPNRWATSI29OCWST@cc.usu.edu> for Symanzik@sunfs.math.usu.edu;
 Sun, 21 May 2000 11:25:12 MDT
Date: Sun, 21 May 2000 11:25:12 -0600
From: weipingdeng <weipingdeng@cc.usu.edu>
Subject: Homework 6 (Part 3)
Sender: weipingdeng <weipingdeng@cc.usu.edu>
To: Symanzik@sunfs.math.usu.edu
Message-id: <3926E03A@webmail.usu.edu>
MIME-version: 1.0
X-Mailer: WebMail (Hydra) SMTP v3.61
Content-type: text/plain; charset="US-ASCII"
Content-transfer-encoding: 7bit
X-WebMail-UserID: weipingdeng
X-EXP32-SerialNo: 00002751
Content-Length: 12066
Status: RO
X-Mozilla-Status: 8001
X-Mozilla-Status2: 00000000
X-UIDL: 38baa60a000002fb

Why did I use the teaching tool at http://cyberk.com/1.0, when I gave the 
lecture on the topic Simple Linear Regression.


The World Wide Web(WWW) has opened up a new medium for education. There are a 
lot of WWW locations that provide various forms of information for statistics 
education in electronic form. For example, electronic textbooks, teachware 
tools and other course materials are availabe at some WWW sites.

In most of the electronic forms, hypertext documents are designed for user 
interaction. In the context of learning statistics, the student can view 
demos, cut and paste sample command into the program, try out suggestions, 
quickly find the definitions for new words, and follow up on related concepts. 
For example, in HTML, help for concepts and keywords can be linked inside the 
document, so that simply clicking on the word allows the user to easily find 
definitions for unfamiliar concepts.

The most impressive thing is dynamic graph which is siginificantly different 
from the classic textbook and teaching method. Using Lisp-stat and Java, a 
series of dynamic graphics have been created to help students to learn 
introductory statistical concepts. The goal of these Web-based demonstrations 
is to present statistical ideas in vivid and direct manner. This is done by 
visualizing concepts that are difficult to communicate using dtraditional 
methods and by eliminating extraneous steps. Adding the interactive components 
would stimulate student interests by providing a 'hands on' learning 
experience. Many topics can be presented using colorful graphics that can be 
manipulated in real-time. The dynamic Web environment provides a powerful tool 
for introductory statistics courses.

After reading some papers on Web-based teaching, referring electronic 
textbooks and teaching tools given in class, and considering the topic Linear 
Regression, I experimented with several tools and textbooks before I decided 
to choose which one as a tool for my lecture.

1. I tried Web-enhanced introductory statistics courses at West Virginia at 
http://www.stat.wvu.edu/srs. A series of Web-based statistical modules and 
infrastructure component are available. Each module covers a statistical topic 
and contains links to images, examples, applets and exercises. If lectures 
notes are simply transferred to the Web as static pages, the gain will be 
minimal. West Virginia University created a dynamic Web environment for 
teaching. I chose the topic regression in the inferential statistical module. 
In the basic principle, a plot of some points are shown. If I click on the 
model
for regression, then the regression line is drawn immediately. If I drag some 
points to different places, the changed regression line will be given too. 
This applet can let students do the interactivity by adding or changing the 
points by themselves and then see how the regression line change. This 
advantage can help students to visualize how the points affect the regression 
line. But this page does not explain the reason and the criterion for drawing 
this regression line, just giving the regression line automatically once the 
points are placed. I would like to show why this kind of line is chosen as the 
regression line, so I did not use this applets in my lecture.


2.I located an applet at 
http://www.stat.sc.edu/~west/javahtml/Regression.html. This applet is similar 
to the above one. It is designed to teach students the effects of leverage 
points on a regression line. Students may add points to the plot by clicking 
the mouse button. The different thing is that it has a summary on how the 
points affect the regression analysis. This should be more helpful. As I said 
in the part one, there is no criterion explained about the regression line, so 
I did not use it in my lecture.


3. I experienced with Hyperstat electronic textbook at 
http://davidmlane.com/hyperstat.  It contains many contents. There are 18 
chapters. All the chapters have the same organization and features. I clicked 
the chapter 15, prediction. Just like our classic textbooks, the chapter was 
divided into several sections. The contents are discussed step by step in 
detail. It is well organized. In the section 1, the text is complete in 
explaining the basic concepts. As I read the text, it seems that I was reading 
a printed textbook except that some description of concepts is accessible 
easily by just clicking the key words. From this text I can understand what 
the regression line is, how the related estimates are computed, and what is 
the criterion. Inside Hyperstat textbook, the text is not significantly 
different from what one encounters in a classical printed textbook. While such 
information is certainly useful and easily accessible, it is static in that it 
cannot respond to user input.

It is helpful that there are a lot of links of references related to the 
chapter, including Analysis tools, Instructional Demos, and Text. I am 
interested in the Demos, so I clicked some of them. In the Demo Regression by 
Eye, if I click New Data button and Draw regression line button, the different 
regression lines will appear with the different sets of points in the picture. 
At the same time, the correlation and the minimun MSE are shown. This Demo 
helps students to visualize how the spread shape of the points determine the 
correlation and the regression line. In the Demo Linear Regression, there is 
an applet that let us mark and change the locations of some ponts, and then 
the equation and graph of the regresson line, and residuals will be given 
accordingly. Of course, these interactivity will be more effective and 
impressive than the classic static textbook.

Since these links are outside the hyper textbook and sometimes it can not be 
open, it is a little bit inconvenient. I prefer the interactivities and 
demonstrations are within the text, so I didn't use it in my lecture.

4. I located a statistical applet at 
http://ww3.whfreeman.com/test/statistics/CorrelationRegression.html. The 
definition of the Least-Square regression line is given first. In the 
following picture, if we add any point and click the least-square line button, 
the regression line will be shown immediately. The funny thing is that we can 
draw any line in the picture and later we can put it in the trash. Students 
may draw the guessed regression line according to the points, and then compare 
it with the accurate regression line. The Mean X and Mean Y lines are also 
available. I think this applet is too simple to explain everything clear about 
Simple Linear Regression.

5. I explored the Berkeley electronic textbook at 
http://www.stat.berkeley.edu/users/stark/sticiGui/text/index.htm. In this text 
index, Home, text, Glossary, problems, calculator, tools, review, and grades 
are availale. Under the Glossary, there are 27 chapters. It contain so many 
topics and contents that it is useful at the undergraduate, graduate and even 
higher level. I chose the chapter 5 Regression, and went to the part of 
Regression line. In the picture for the example, both the SD line and 
regression line are drawn. That help students to compare the two kinds of 
lines. The data set and the unvariate statistics are also available. Following 
this picture, some exercises are given. Students can get the feedback 
immediately after typing the answer. This is pretty good. Then, the detailed 
explanation about the related concepts, the criterion, estimate computations 
are given clearly. Some following exercises are given again to check if 
students understand the concepts and computations well, and the solutions are 
attached. The special cases of the regression line are dicussed further, 
followed by the examples, exercises and solutions. This part is complete for 
simple linear regression.

I would like to watch some dynamic graph, so I clicked the tools in the index 
and located at an applet for Correlation and Regression. We can use the slider 
to change the correlation and the number of data. The regression line and the 
residual plot can be shown. Points can be added or cleared. It is fun and 
helpful to play in this applet.

Because the text for simple linear regression is a little more complicated in 
this textbook, so I didn't choose it for my introductory statistical lecture.

6. I located the sufstat textbook at 
http://surfstat.newcastle.edu.au/surfstat/mail/surfstat.html. All the material 
I can find inside the textbook about the linear regression is the text just as 
in the classic printed textbook. There is no dynamic graph available, so I 
didn't use it in my lecture.

7. I chose the teaching tool at http://cyberk.com/1.0/D-2/indes.html, in my 
lecture. In each unit, there are Think First, Three Keys, Practice, Resources, 
Feedback and Calculator. The reason why I chose this teaching tool is that it 
has a lot of interactivities and dynamic graphics in examples and exercises. 
Since I only have 18 minutes to talk about the Simple Linear Regression, I 
just used the Basics for my lecture. In this page, the linear model equation 
is written. How to fit a best straight line if the data are given? What 
criterion can we use? In order to talk about these problems, I used the first 
interactivity example. In this picture, the points are marked already. Our 
goal is to select a straight line that best fit the data. Two sliders can be 
moved to adjust the straight line. Following the picture, X & Y, a+bX, 
Residuals and Square of the Residuals are listed in the table. The Sum of 
Square of the Residuals (SSR), Estimate of intercept a, and Estimate of slope 
b are shown on the top of the picture. I explain the relationship between the 
value in the table and the corresponding coordinate of the points or length of 
the line segment. The different color for the points and lines make everything 
clear. The best line should get as close to the points as possible, so the 
absolute values of the residuals (the lengths of the vertical red line 
segments) should be as small as possible. Considering all the residuals, one 
criterion (Least Square) comes out. It is to select the straight line such 
that the SSR is smallest. Then according to the value of the SSR on the top of 
the picture to adjust the line untill the SSR is minimum. This line is called 
least square line, and the estimates are call Least Square estimates. I think 
this interactivity is impressive and effective, because it demonstrate how the 
regression line is obtained vividly. Then one exercise was practiced by fellow 
students in class to find the least square line, with slider to adjust the 
lines to get the smallest SSR. The feedback is given immediately. Using this 
electronic teaching tool is absolutely much more effective than the classic 
pure 'chalk and textbook' method. In this lecture, students can visualize the 
whole dynamic process of the graph and play it by themselves. I just talked 
the Basics 1 and 2 in detail. Actually, Basics 3 gives the formulas to compute 
the estimates of the regression line by the Least square criterion. If I had 
more time, I would compare these two methods.

These are the reasons why I chose the http://Cyberk.com/1.0 but not the 
others.


Rreferences:

A.J.Rossini and Rosenberger. (1994), Teaching Statistics and Computing via 
Multimedia through the Would Wide Web, Feature Article, Statistical Computing 
& Graphics, vol.5 No.3, 10-13.

R.Webster West and R. Todd Ogdent. (1998), Interactive Demonstrations for 
Statistics Education on the World Wide Web. Journal of Statistics Education 
v.6,n.3. http://www.amstat.org/publications/jse/v6n3/west.html

Jan de Leeuw. (1997), The UCLA Statistics Textbook and Modules, Isi Bulletin, 
Book 2, 55-58.

E. James Harner and William C. Wojciechowski. (1998), A Web-enhanced 
Introductory Statistics Course, Computing Science and Statistics, 
V.29,n.1,316-326.

Robin H. Lock. (1997), Internet Resources for Teaching Statistics, Computing 
Science and Statistics, V.29,N.2, 339-343.

